\section{Specifications and Requirements}
The implementation of the multi-task deep learning framework for cattle health monitoring requires a specific configuration of hardware and software components to ensure real-time inference and model training.

\subsection{Hardware Requirements}
\begin{itemize}
    \item \textbf{Development and Training Unit:} The model training was executed on a high-performance research workstation equipped with an NVIDIA GeForce RTX 4080 GPU (16GB VRAM, Ada Lovelace architecture supporting FP8 and FP16 tensor operations), an Intel Core i9-13900K CPU, 64GB of DDR5 RAM, and a 2TB NVMe SSD to handle high-throughput image and video loading.
    \item \textbf{Edge Deployment Unit:} For live farm surveillance, the system is designed to run on resource-constrained edge computers such as the Jetson Orin Nano (8GB VRAM) or a Raspberry Pi 5. The total model size is optimized to remain under 10 million parameters ($<$ 40MB file size in FP16 format) to fit into edge device memory. Furthermore, the shared EfficientNet-B0 backbone requires approximately 0.39 GFLOPs per frame. Processing a full 20-frame spatiotemporal sequence requires only $\sim$7.8 GFLOPs, which falls comfortably within the real-time processing capabilities of a Jetson Orin Nano (capable of 40 TOPS).
    \item \textbf{Imaging Sensors:} Standard 1080p surveillance cameras operating at 30 frames per second (FPS) positioned at cattle walkways or feeding troughs. For the Dryad dataset, depth maps were captured using active infrared depth sensors.
\end{itemize}

\subsection{Software Requirements}
\begin{itemize}
    \item \textbf{Operating System:} Windows 11 / Linux (Ubuntu 22.04 LTS).
    \item \textbf{Programming Language:} Python 3.12.3.
    \item \textbf{Deep Learning Library:} PyTorch 2.5.1 + CUDA 12.1.
    \item \textbf{Backbone Repository:} PyTorch Image Models (\texttt{timm}) library for loading pre-trained EfficientNet-B0 and MobileNetV3 architectures.
    \item \textbf{Computer Vision Tools:} Computer vision operations are implemented using OpenCV 4.9.0 for video processing and Albumentations 1.4.0 for image augmentation and preprocessing tasks.
    \item \textbf{Data Pipelines:} Pandas, NumPy, and Scikit-learn for group-wise dataset splits and metric computations.
\end{itemize}

\section{Societal Impact}
The development of automated cattle health monitoring systems has a profound positive impact on the agricultural sector and animal welfare:
\begin{itemize}
    \item \textbf{Enhancement of Animal Welfare:} Early detection of lameness and abnormal behaviors (such as reduced eating or drinking) allows farmers to treat sick cows days before clinical symptoms become visible, minimizing animal suffering.
    \item \textbf{Support for Precision Livestock Farming (PLF):} Transitioning from manual, subjective herd inspections to continuous 24/7 automated monitoring reduces human error and standardizes assessments such as Body Condition Scoring.
    \item \textbf{Agricultural Labor Relief:} As intensive dairy farming expands, the ratio of cows per stockperson increases. Automated monitoring alleviates labor shortages, allowing farm workers to focus on specialized veterinary interventions rather than manual surveillance.
\end{itemize}

\section{Environmental Impact}
Deep learning models are historically compute-intensive, contributing to carbon emissions during training. Our proposed framework addresses this environmental concern through:
\begin{itemize}
    \item \textbf{Green Computing via Multi-Task Learning (MTL):} Instead of deploying four separate convolutional neural networks (each consuming GPU power and memory), our hard parameter sharing approach utilizes a single shared EfficientNet-B0 encoder. This consolidation is expected to reduce duplicated backbone computation compared with four separate models, although measured edge-device energy benchmarking remains future work, minimizing the operational carbon footprint of edge surveillance systems deployed on thousands of farms.
    \item \textbf{Resource Optimization:} By implementing sparse temporal sampling (reducing 300-frame video clips to 20 key frames) and capping dataset classes, we decreased active training GPU runtimes, reducing energy overhead during the development lifecycle.
\end{itemize}

\section{Ethical Issues}
The use of visual monitoring systems on farms raises ethical concerns that must be addressed:
\begin{itemize}
    \item \textbf{Non-Invasive Monitoring:} Traditional cattle identification methods such as branding, ear tags, and implanted RFID devices can cause pain and can negatively affect animal welfare. Our proposed vision-based monitoring approach avoids physical contact and reduces stress for the animals.
    \item \textbf{Data Privacy:} Local edge processing helps prevent farm data and operational details from being exposed to public networks. The proposed offline edge system ensures data privacy and security.
\end{itemize}

\section{Standards - if applicable}
The proposed system and agricultural metrics are developed according to the following standards:
\begin{itemize}
    \item \textbf{Body Condition Scoring (BCS) Standards:} The Body Condition Scoring system uses globally recognized 5-point Elanco scale with 0.25 increments in ScienceDB and integer scores in Dryad.
    \item \textbf{Locomotion Scoring Standards:} Lameness labels are based on the Sprecher scoring system, where cattle are classified as either healthy or lame.
    \item \textbf{IEEE Software Quality Standards:} The software is developed according to the PEP 8 Python style guide, and the evaluation metrics use standard definitions such as Macro F1-score, MAE, and AUC-ROC.
\end{itemize}

\section{Project Management Plan}
The workload was divided among five team members according to their base backbone assignments and particular tasks to handle the multi-task and multi-backbone complexity of the project:
\begin{itemize}
    \item \textbf{Hasin Ishrak:} Preprocessed the main datasets, implemented the baseline training pipelines, evaluated the ResNet-18 baseline, and executed the primary ablation studies including CBAM evaluations and the comparison between RGB and Depth modalities. 
    \item \textbf{Namira Abrar Haque:} Evaluated the MobileNetV3-Small baseline and examined how highly compressed networks perform for resource-constrained edge deployment. 
    \item \textbf{Sanjida Akter Bithi:} Evaluated the ResNet-50 baseline, and evaluated the advantages of deeper architectures on spatial feature extraction. 
    \item \textbf{Shouvik Banik:} Evaluated the DenseNet121 baseline, and studied the effect of dense feature connectivity in multi-class behavior recognition. 
    \item \textbf{Nusrat Lamya Faruk:} Evaluated the winning EfficientNet-B0 baseline, developed the final spatiotemporal multi-task model, combined the LSTM sequence-tracking layers, and implemented the real-time visualization interface. 
\end{itemize}
The timeline was structured in sprints: dataset curation, single-task baseline training, backbone selection, spatiotemporal multi-task model integration, and finally report writing and visualizer deployment.

\section{Risk Management}
\begin{table}[H]
\centering
\caption{Project Risks and Mitigation Strategies}
\label{tab:risk_management}
\resizebox{\textwidth}{!}{%
\begin{tabular}{l>{\raggedright\arraybackslash}p{5cm}>{\raggedright\arraybackslash}p{6cm}}
\toprule
\textbf{Risk Identified} & \textbf{Potential Impact} & \textbf{Mitigation Strategy} \\
\midrule
Class Imbalance & Model fails to learn rare behaviors (licking, drinking). & Implemented Focal Loss ($\gamma=2$) and hard data capping (max 3000 images per class). \\
Data Leakage & Memorization of cow identities leading to overfitting. & Task-appropriate splitting (cow-wise for BCS/behavior, video-wise for lameness, image-wise within identity for ID). \\
Out of Memory (OOM) & GPU VRAM crash when processing long sequences. & Integrated Automatic Mixed Precision (AMP) in FP16 and sparse temporal sampling of 20 frames. \\
Domain Shift & Performance drop when deployed on unseen farms. & Performed cross-dataset validation on CBVD-5 to evaluate generalizability and recommend future unfreezed fine-tuning. \\
\bottomrule
\end{tabular}%
}
\end{table}

\section{Economic Analysis}
The financial viability of our multi-task framework is compared against traditional farm monitoring systems in Table~\ref{tab:economic_analysis}.
\begin{table}[H]
\centering
\caption{Economic Cost-Benefit Analysis}
\label{tab:economic_analysis}
\resizebox{\textwidth}{!}{%
\begin{tabular}{l>{\raggedright\arraybackslash}p{6cm}>{\raggedright\arraybackslash}p{6cm}}
\toprule
\textbf{Component} & \textbf{Traditional Sensor Systems} & \textbf{Proposed MTL Edge System} \\
\midrule
\textbf{Sensor Hardware} & Wearable collars, step counters, and RFID tags (\$80 - \$120 per cow). For 500 cows: \$40,000+. & Standard IP CCTV cameras (\$150 each) + edge computing unit (\$300). Total: \$1,500. \\
\textbf{Maintenance} & Battery replacements, collar damage, and physical tag losses (15\% annual loss). & Minimal hardware contact, camera cleaning and software updates only. \\
\textbf{Computational Cost} & High cloud subscription fees for running separate analysis pipelines. & Offline local edge inference with zero recurring cloud subscription costs. \\
\bottomrule
\end{tabular}%
}
\end{table}
The analysis demonstrates that our visual multi-task edge framework is highly scalable, lowering the initial investment barrier for small-to-medium scale dairy farms by over 90\%.